\documentclass[12pt,a4paper]{article}
\usepackage[utf8]{inputenc}
\usepackage[T1]{fontenc}
\usepackage{amsmath,amssymb,amsthm}
\usepackage{geometry}
\usepackage{hyperref}
\usepackage{booktabs}

\geometry{margin=1in}

\newtheorem{theorem}{Theorem}
\newtheorem{axiom}{Axiom}
\newtheorem{proposition}{Proposition}
\newtheorem{lemma}{Lemma}

\title{\textbf{Fixed-point existence of self-consistent semiclassical gravity}\\[6pt]
{\large DRAFT}}

\author{Will Regelmann\thanks{Corresponding author.} \quad Claude (Anthropic)\thanks{AI
assistant; contributed to mathematical development and analysis.}}

\date{February 2026}

\begin{document}
\maketitle

\begin{abstract}
The semiclassical Einstein equation $G_{\mu\nu}[g] = 8\pi G\langle\hat{T}_{\mu\nu}\rangle_g$
is a fixed-point condition: the geometry and quantum state must be mutually consistent.
We address the question of whether such fixed points exist. We present three analyses:
(1)~exact existence in the cosmological case via the trace anomaly
(Starobinsky, 1980) (existence \textbf{Rigorous}, by citation; the explicit fixed-point
curvature scale $H_0^2 = 180\pi/(G|a_2|)$ is \textbf{Sketch} --- a convention-dependent
reduction not verbatim in the cited sources, Section~\ref{sec:starobinsky}); (2)~a perturbative argument near
classical solutions \emph{on globally hyperbolic spacetimes with compact Cauchy surfaces},
via a Banach contraction whose target estimate is a contraction constant
$\kappa \sim (m/M_P)^2 \ll 1$ for massive fields, developed via a Hadamard decomposition
of the order-reduced response kernel into a local differential operator and a non-local
smooth kernel --- this argument is presently a \textbf{Sketch}: an earlier Rigorous label
was withdrawn after review identified gaps in the kernel bound and in the contraction
estimate (Section~\ref{sec:banach}), and the massless case, relevant to photons and gluons,
remains open (Section~\ref{sec:open}); and (3)~conditional existence in the general
case via the Schauder fixed-point theorem on globally hyperbolic spacetimes with compact
Cauchy surfaces, under six stated assumptions of which one remains an open conjecture
(the proof as written also has identified repairable gaps; Section~\ref{sec:schauder}).
An earlier claim that the Planck scale emerges from these analyses as the natural validity
boundary is withdrawn: the derivation chain behind the contraction constant contains a
dimensional inconsistency, and neither of its two candidate readings places the breakdown
at the Planck scale (Section~\ref{sec:planck}). An effective initial-value formulation,
valid order-by-order in $\kappa$, is derived from the global fixed-point condition.

\bigskip\noindent\textit{Note:} This is a companion paper to ``Gaussian temporal profile
of gravitational decoherence from the Einstein-Langevin equation.'' The decoherence
prediction of that paper does not depend on the existence results derived here; the present
paper addresses the independent question of whether the semiclassical framework is
internally self-consistent.
\end{abstract}

%======================================================================
\section{Introduction}
%======================================================================

The semiclassical Einstein equation:
\begin{equation}
G_{\mu\nu}[g] = 8\pi G\,\langle\Psi|\hat{T}_{\mu\nu}[g]|\Psi\rangle
\label{eq:sce}
\end{equation}
is not a dynamical equation in the usual sense. It is a fixed-point condition: we seek the
geometry $g_{\mu\nu}$ such that the quantum state it hosts, evolving on it, sources exactly
that geometry. The position that eq.~\eqref{eq:sce} may be fundamental (rather than an
approximation to quantized gravity) dates to M\o ller (1962) and Rosenfeld (1963) and has
been developed by Oppenheim~\cite{oppenheim}. Whether this position is internally
consistent depends on whether fixed points exist.

This paper presents existence arguments at three levels of generality. The cosmological
result (Section~\ref{sec:starobinsky}) is exact. The perturbative Banach-contraction
argument (Section~\ref{sec:banach}) is labeled \textbf{Sketch}: it was promoted to
Rigorous in an earlier revision, and that label has been withdrawn after review identified
specific gaps, catalogued where they occur. The Schauder argument
(Section~\ref{sec:schauder}) remains conditional on an open conjecture~(A3), with
additional repairable gaps noted in its proof. An earlier claim that the Planck scale
emerges from the convergence analysis as the natural validity boundary is withdrawn
pending a corrected contraction estimate (Section~\ref{sec:planck}).

\medskip\noindent The perturbative contraction estimate is permanently Sketch at the stated generality: the binding obstruction (Gap~M3) is mass-independent --- the non-local response kernel is a Lorentzian Green operator with causal-past support, not a spatial elliptic resolvent, and this distinction does not depend on field mass --- so the same Gap~(M3) obstruction applies equally to massless fields such as the photon and gluon (Section~\ref{sec:open}). The estimate also requires a compact Cauchy surface, so it does not cover Minkowski space or asymptotically flat backgrounds (Section~\ref{sec:open}).

%======================================================================
\section{The self-consistency map}
\label{sec:map}
%======================================================================

Define $\mathcal{F}: g_{\mu\nu} \mapsto g'_{\mu\nu}$ by:
\begin{enumerate}
\item Given trial geometry $g$, construct the QFT on $(\mathcal{M},g)$ and identify the
physical state $|\Psi_g\rangle$ satisfying the Hadamard condition.
\item Compute $S_{\mu\nu}[g] \equiv 8\pi G\langle\hat{T}_{\mu\nu}\rangle_{g,\text{ren}}$.
\item Solve $G_{\mu\nu}[g'] = S_{\mu\nu}[g]$ for $g'$.
\end{enumerate}
We seek $g^*$ such that $\mathcal{F}(g^*) = g^*$.

\medskip\noindent\textbf{State selection.} Step~1 requires care: the Hadamard condition
constrains the wavefront set of the two-point function but does not uniquely select a
state. In the cosmological case (Section~\ref{sec:starobinsky}), de~Sitter symmetry
selects the conformal vacuum. In the perturbative case (Section~\ref{sec:banach}), the
state on the corrected geometry is the unique perturbation of the state on the classical
background. For the Schauder argument (Section~\ref{sec:schauder}), multi-valuedness is
handled by restricting to a suitable branch or by applying Kakutani's theorem for
set-valued maps.

\medskip\noindent\textbf{Single-valuedness assumption.} The Banach argument (Section~\ref{sec:banach}) assumes $\mathcal{F}$ is single-valued: in the perturbative regime, the Hadamard state on the corrected geometry is the unique perturbation of the state on the classical background, so no state-selection ambiguity arises. The Schauder argument (Section~\ref{sec:schauder}) likewise assumes single-valuedness via the branch restriction noted above. An extension to the genuinely set-valued case via Kakutani's fixed-point theorem is possible but is not pursued here.

%======================================================================
\section{Cosmological fixed point: exact solution}
\label{sec:starobinsky}
%======================================================================

\begin{theorem}[Starobinsky, 1980 \cite{starobinsky}]
The self-consistency equation on FRW spacetime with conformal matter admits an exact
de~Sitter solution $a(t) = a_0 e^{H_0 t}$ with:
\begin{equation}
H_0^2 = \frac{180\pi}{G|a_2|}
\end{equation}
where $a_2$ is the coefficient of the Euler density in the trace anomaly. The de~Sitter
solution is \emph{unstable}: the quasi-de~Sitter stage is long-lived but decays, and in
the inflationary application this instability is the graceful-exit mechanism.

\emph{(Existence and instability of the de~Sitter fixed point: \textbf{Rigorous}, by
citation to Starobinsky~\cite{starobinsky}. Explicit coefficient
$H_0^2 = 180\pi/(G|a_2|)$: \textbf{Sketch} --- see Rigor status note below.)}
\end{theorem}

\noindent\emph{Proof sketch.} For constant $H$, the trace anomaly~\cite{capper_duff}
reduces to $\langle\hat{T}^\mu{}_\mu\rangle = -a_2 H_0^4/120\pi^2$. The Friedmann
equation yields $H_0$. The instability of the de~Sitter stage follows from the linearized
equation for $\delta H$~\cite{starobinsky}. \hfill$\square$

\smallskip\noindent\textbf{Rigor status (red-team demotion, June 2026).} The
\emph{existence and instability} of the anomaly-driven de~Sitter fixed point are
\textbf{Rigorous, by citation} to Starobinsky~\cite{starobinsky}. The \emph{explicit
coefficient} $H_0^2 = 180\pi/(G|a_2|)$ is \textbf{Sketch}: the algebra is internally
consistent (combining $\langle\hat{T}^\mu{}_\mu\rangle = -a_2 H_0^4/120\pi^2$ with the
trace of the semiclassical equation, $-R = 8\pi G\langle\hat{T}^\mu{}_\mu\rangle$ and
$R = 12H_0^2$ on de~Sitter, gives $H_0^2 = 180\pi/(Ga_2)$), but two steps are not closed
by the cited sources. (i)~Starobinsky parametrizes the consistency condition through the
effective action of $N$ conformal fields (a curvature-scale bound set by the field content),
not as the closed form $180\pi/(G|a_2|)$ with ``$a_2 = $ Euler-density coefficient'';
the $180\pi$ is a later reparametrization. (ii)~Capper--Duff~\cite{capper_duff} establishes
the \emph{existence} of the conformal trace anomaly, but the de~Sitter-evaluated,
field-content-summed normalization $\langle\hat{T}^\mu{}_\mu\rangle = -a_2 H_0^4/120\pi^2$
requires fixing a field content and an anomaly-coefficient convention, neither of which that
source supplies. Relative to the standard split $\langle\hat{T}^\mu{}_\mu\rangle =
c\,C^2 - a\,E$ with $C^2 = 0$ and $E = 24H_0^4$ on de~Sitter, the factor $120\pi^2$ amounts
to $a_2 = 2880\pi^2\,a$ --- a convention internal to this paper. Finally, the absolute value
in $|a_2|$ silently encodes a positivity condition: $H_0^2 > 0$ requires $a_2 > 0$ (which
holds for standard conformal matter, where the type-A coefficient $a > 0$, but is not free).
Issue~\#133 tracks restoring this to Rigorous (fix the convention, cite a modern
anomaly-induced-inflation reference for the explicit coefficient, state the $a_2 > 0$
condition).

\smallskip\noindent\emph{Remark (correction).} An earlier draft of this paper asserted
that the de~Sitter solution is a stable attractor with exponentially damped perturbations.
This misstates the cited source: in Starobinsky's model the de~Sitter stage is unstable
(its decay is what ends inflation), and the existence claim --- not stability --- is what
this section uses. Only existence of the exact fixed point is relied upon below.

This constitutes an existence proof: the de~Sitter geometry, used to compute the quantum
stress-energy of fields living on it, sources exactly itself.

%======================================================================
\section{Perturbative fixed point: Banach contraction}
\label{sec:banach}
%======================================================================

\noindent\textbf{Rigor status: Sketch (demoted from Rigorous).} The results of this
section --- Lemma~\ref{lem:kernel_bound} and Theorem~\ref{thm:banach} --- were promoted
to Rigorous in an earlier revision (via the Hadamard-decomposition argument recorded in
\texttt{explorations/2026-03-02-A1-Hs-bound-derivation.md}). That label has been
withdrawn: review identified four gaps, labeled (M1), (M3), (M4), (M5) below, that the
arguments do not close. The proofs are retained as the record of the sketch; each gap is
flagged where it occurs. Closing them is follow-up work (Section~\ref{sec:open}).

\medskip

Let $\bar{g}$ be a solution of the classical Einstein equations. Define the
quantum-corrected map:
\begin{equation}
\mathcal{F}(g) = \bar{g} + (G^{\bar{g}}_{\text{lin}})^{-1}\!\left[8\pi G\langle\hat{T}_{\mu\nu}\rangle^{(\text{quantum})}_g\right]
\end{equation}

\noindent We first establish a bound on the order-reduced response kernel, then use it
to prove the contraction estimate.

\begin{lemma}[Order-reduced response kernel bound --- \textbf{Sketch}]
\label{lem:kernel_bound}
Let $\phi$ be a massive scalar field with mass $m > 0$ on a globally hyperbolic spacetime
$(\mathcal{M},\bar{g})$ with compact Cauchy surface~$\Sigma$. After Parker--Simon order
reduction~\emph{\cite{parker_simon}}, the order-reduced response kernel
$\mathcal{K}^{\mathrm{red}}_{\mu\nu\alpha\beta}(x,x')
\equiv \delta\langle\hat{T}_{\mu\nu}\rangle^{\mathrm{red}}_{\mathrm{ren}}/\delta g^{\alpha\beta}(x')$
satisfies
\begin{equation}
\bigl\|\mathcal{K}^{\mathrm{red}}\bigr\|_{H^s \to H^{s-2}}
\;\leq\; C_{\mathcal{K}} \cdot \frac{\hbar\, m^2}{(4\pi)^2}
\label{eq:kernel_bound}
\end{equation}
uniformly for all $g \in K_\rho$, where $C_{\mathcal{K}}$ is a dimensionless constant
depending on $\rho$, $s$, and the background geometry~$\bar{g}$.
\end{lemma}

\noindent\emph{Proof.} We decompose the response kernel via the Hadamard
parametrix into a local differential operator and a non-local smooth kernel,
bound each separately, and combine via the triangle inequality.

\smallskip\noindent\textit{Step~1: Hadamard decomposition of the response kernel.}
By assumption~(A1), the renormalized stress-energy tensor is obtained by
point-splitting with Hadamard subtraction~\cite{hollands_wald_wick}:
\begin{equation}
\langle\hat{T}_{\mu\nu}(x)\rangle_{\mathrm{ren}}
= \lim_{x'\to x}\mathcal{D}_{\mu\nu}(x,x')\,W(x,x') + C_{\mu\nu}[g](x)
\label{eq:hadamard_ren}
\end{equation}
where $W = W_2 - H$ is the smooth Hadamard remainder ($W_2$ the full
two-point function, $H$ the universal Hadamard singular part determined by
local geometry via the Hadamard recursion), $\mathcal{D}_{\mu\nu}$ is the
second-order point-splitting differential operator, and $C_{\mu\nu}[g]$ collects
the finite local curvature counterterms (involving $R$,
$R_{\mu\alpha\nu\beta}R^{\alpha\beta}$, $\Box R$).

The functional derivative $\delta/\delta g^{\alpha\beta}(y)$ of
eq.~\eqref{eq:hadamard_ren} yields three types of contributions to the
response kernel $\mathcal{K}_{\mu\nu\alpha\beta}(x,y)$:
\begin{enumerate}
\item[(I)] Varying the smooth remainder~$W$ --- non-local, mediated by the
massive field propagator;
\item[(II)] Varying the point-splitting operator $\mathcal{D}_{\mu\nu}$ ---
local, from Christoffel symbol variation at $y=x$;
\item[(III)] Varying the curvature counterterms $C_{\mu\nu}$ --- local.
\end{enumerate}

\smallskip\noindent\textit{Step~2: Order reduction and local/non-local decomposition.}
Before order reduction, Term~(III) involves up to four derivatives of the
metric perturbation (from the $\Box R$ counterterm --- the source of the
Horowitz runaway instability~\cite{horowitz}). After Parker--Simon order
reduction~\cite{parker_simon} (assumption~(A5)), the fourth-derivative
terms are perturbatively eliminated using the classical Einstein equation,
leaving Term~(III) as a differential operator of order~$\leq 2$.

Term~(II) is a differential operator of order~$\leq 1$: varying
covariant derivatives introduces at most one derivative of~$\delta g$
via the Christoffel symbols. Term~(I) decomposes into contact
contributions at $y = x$ (absorbed into the local part) and a smooth
non-local kernel for $y\neq x$.

Define $\mathcal{K}^{\mathrm{red}}_{\mathrm{loc}}$ as the sum of all
local (distributional) contributions: a differential operator of
order~$\leq 2$ on~$\Sigma$. Define $\mathcal{K}_{\mathrm{nl}}$ as the
remaining non-local smooth kernel.

\smallskip\noindent\textit{Step~3: Bound on the local part.}
The local part takes the form
$\mathcal{K}^{\mathrm{red}}_{\mathrm{loc}}
= \sum_{|\gamma|\leq 2}a_\gamma(x)\,\partial^\gamma$
with smooth coefficients whose $L^\infty(\Sigma)$ norms are set by the
one-loop scale: $\|a_\gamma\|_{L^\infty} \leq C_{\mathrm{loc}}
\cdot\hbar m^2/(4\pi)^2$. A differential operator of order~$d$ with
$L^\infty$ coefficients maps $H^s(\Sigma)\to H^{s-d}(\Sigma)$
boundedly (by the Leibniz rule and Sobolev embedding on compact
manifolds). Therefore:
\begin{equation}
\bigl\|\mathcal{K}^{\mathrm{red}}_{\mathrm{loc}}\bigr\|_{H^s\to H^{s-2}}
\;\leq\; C_{\mathrm{loc}}\cdot\frac{\hbar\,m^2}{(4\pi)^2}.
\label{eq:loc_bound}
\end{equation}

\smallskip\noindent\textit{Step~4: Bound on the non-local part.}
For a massive field ($m > 0$), the smooth remainder~$W$ depends on the
metric through the Klein--Gordon equation
$(\Box_g + m^2 + \xi R)\phi = 0$. Its functional derivative
$\delta W/\delta g^{\alpha\beta}(y)$ is mediated by the massive
Klein--Gordon resolvent on~$\Sigma$, whose integral kernel decays
exponentially: $|(-\Delta_\Sigma + m^2 + \xi R|_\Sigma)^{-1}(x,y)|
\leq C\,e^{-m\cdot d_\Sigma(x,y)}$ for $d_\Sigma(x,y) \gg 1/m$. This
is a standard property of elliptic resolvents with positive mass on
compact Riemannian manifolds.

Consequently, $\mathcal{K}_{\mathrm{nl}}$ has a smooth kernel on
$\Sigma\times\Sigma$ with exponential decay. A smooth integral operator
on a compact manifold is a smoothing operator, bounded
$H^s\to H^{s'}$ for all~$s,s'$; its $H^s\to H^{s-2}$ operator norm
is bounded by its Hilbert--Schmidt norm:
\begin{equation}
\bigl\|\mathcal{K}_{\mathrm{nl}}\bigr\|_{H^s\to H^{s-2}}
\;\leq\; \|K_{\mathrm{nl}}\|_{L^2(\Sigma\times\Sigma)}
\;\leq\; C_{\mathrm{nl}}\cdot\frac{\hbar\,m^2}{(4\pi)^2}.
\label{eq:nl_bound}
\end{equation}

\smallskip\noindent\textit{Step~5: Combined bound.}
By the triangle inequality,
$\|\mathcal{K}^{\mathrm{red}}\|_{H^s\to H^{s-2}}
\leq \|\mathcal{K}^{\mathrm{red}}_{\mathrm{loc}}\|
+ \|\mathcal{K}_{\mathrm{nl}}\|
\leq C_{\mathcal{K}}\cdot\hbar m^2/(4\pi)^2$,
establishing~\eqref{eq:kernel_bound}.

\smallskip\noindent\textit{Step~6: Uniformity over $K_\rho$.}
By assumption~(A2), the map $g \mapsto \mathcal{K}^{\mathrm{red}}[g]$ is continuous.
The set $K_\rho$ is compact in $H^{s'}$ for $s' < s$ by Rellich--Kondrachov
(established in Section~\ref{sec:schauder}). A continuous function on a compact set
attains its supremum, so
$\sup_{g \in K_\rho}\|\mathcal{K}^{\mathrm{red}}[g]\|_{\mathrm{op}} < \infty$.
The constant $C_{\mathcal{K}}$ absorbs this supremum. \hfill$\square$

\smallskip\noindent\textbf{Gap (M3): hyperbolic-to-elliptic substitution in Step~4.}
The variation $\delta W/\delta g^{\alpha\beta}(y)$ is governed by the \emph{Lorentzian}
(hyperbolic) Klein--Gordon equation on the spacetime: varying the field equation gives
$(\Box_g + m^2 + \xi R)\,\delta W_2 = -\delta(\Box_g + \xi R)\,W_2$, whose solution
operator has light-cone support. Step~4 instead bounds the variation using the
\emph{spatial elliptic} resolvent $(-\Delta_\Sigma + m^2 + \xi R|_\Sigma)^{-1}$ and its
exponential decay. For massive Lorentzian propagators, exponential decay holds only at
\emph{spacelike} separation; a metric perturbation at $y$ in the causal past of $x$
influences $W(x,x)$ with no exponential suppression. No reduction of the response kernel
to a purely spatial operator on $\Sigma$ is constructed, so the exponential-decay bound
on $\mathcal{K}_{\mathrm{nl}}$ is not established as stated. This is the principal gap.

\emph{This gap is structural, not presentational} (investigated in
\texttt{explorations/2026-06-16-FPE4-M3-structural-obstruction.md}). The non-local
kernel is the image of a source under a Green operator of the \emph{normally hyperbolic}
operator $(\Box_g + m^2 + \xi R)$, whose support is the causal past $J^-(x)$ (for the
bare Green operator; the renormalized-kernel support has not been established to be
smaller); this is an operator-support property, independent of the field mass (the compact-$\Sigma$ spectral
gap $\lambda_1(\Sigma)$ and the conformal shift $\xi R$ control the \emph{spatial}
elliptic operator, which is the wrong operator, so the gap is not cured in the massless
limit either). Reducing this object to an $H^s(\Sigma)\to H^{s-2}(\Sigma)$ operator
presupposes a canonical decomposition of the Lorentzian Green operator into time
evolution off a chosen $\Sigma$ --- i.e.\ a global time function or homogeneous slicing.
Assumptions (A1)--(A6) supply neither (this is the same unjustified background structure
flagged as Gap (M8): a genuine spatial reduction of the kernel would also exhibit the
chosen $\Sigma$ as pure gauge, and conversely). The only known route to a contraction for
the Lorentzian object --- Meda--Pinamonti--Siemssen on FLRW
\cite{meda_pinamonti_siemssen,pinamonti_siemssen} --- contracts in a norm over an interval
of cosmological \emph{time} on the Hubble function, exploiting homogeneous slicing, and is
not a single-Cauchy-surface $H^s$ bound. \textbf{Consequently the perturbative contraction
of this section is permanently Sketch at the stated generality:} a repair requires
geometric input beyond (A1)--(A6) --- a controlled global foliation --- and is then
foliation-restricted, yielding existence ``Rigorous on a [homogeneous / controlled]
slicing,'' local in time, rather than the general statement of Theorem~\ref{thm:banach}.
This does not refute existence of the fixed point, which holds exactly on de~Sitter
(Section~\ref{sec:starobinsky}) and on FLRW \cite{meda_pinamonti_siemssen}; it obstructs only
the single-slice spatial-resolvent \emph{proof strategy} written here.

\smallskip\noindent\textbf{Gap (M4): asserted coefficient scaling in Step~3.}
The bound $\|a_\gamma\|_{L^\infty} \leq C_{\mathrm{loc}}\cdot\hbar m^2/(4\pi)^2$ is
asserted, not derived. It fails for the mass-independent anomaly counterterm responses:
the variation of the $R^2$-, $C^2$-, and Euler-type counterterms in $C_{\mu\nu}[g]$
carries coefficients $\sim \hbar/(4\pi)^2$ times curvature, contributing local terms
scaling as $\hbar R$, not $\hbar m^2$. These can be absorbed into $C_{\mathcal{K}}$ only
at the cost of a factor $R/m^2$, which makes the displayed $m^2$ scaling --- and hence
the $(m/M_P)^2$ headline of Theorem~\ref{thm:banach} --- cosmetic whenever
$R \gtrsim m^2$. A derivation of the coefficient bounds with explicit mass and curvature
dependence is required. This is the same fork the (M1) diagnosis reaches from the
operator-norm side: fixing the single Sobolev reference scale $\mu^2\sim R$ that makes the
$\kappa$ chain consistent forces the local-part scale to $\hbar R/(4\pi)^2$. (M1) and (M4)
must therefore be repaired jointly.

\smallskip\noindent\textit{Remark on the massless case.} The binding
obstruction (Gap~M3) is mass-independent: the non-local response kernel is
a Green operator of the Lorentzian hyperbolic operator $(\Box_g + m^2 + \xi R)$
with causal-past support $J^-(x)$, and no controlled reduction to a spatial
$H^s(\Sigma)$ operator exists within (A1)--(A6) regardless of whether $m > 0$
or $m = 0$. The algebraic-decay and Hilbert--Schmidt arguments that appear in
the massive Sketch are downstream of the spatial-operator assumption that
Gap~(M3) already identifies as illegitimate; they do not constitute an
independent infrared obstruction. The massless case therefore fails for the
same structural reason as the massive case, and remains open on the same terms.

\smallskip\noindent\textit{Remark on the K\"all\'en--Lehmann approach.}
On Minkowski spacetime, the bound can alternatively be established via the
K\"all\'en--Lehmann spectral representation of the stress-energy two-point
function~\cite{phillips_hu,hu_verdaguer}: the spectral gap above $4m^2$
provides polynomial suppression at high momenta, and the Schur test
translates this to an $H^s\to H^{s-2}$ bound. This approach is rigorous on
flat backgrounds but does not extend directly to general curved spacetimes,
where no K\"all\'en--Lehmann representation for the response kernel is
available. The Hadamard decomposition above avoids this difficulty.

\smallskip\noindent\textit{Remark on related work.} Meda, Pinamonti, and
Siemssen~\cite{meda_pinamonti_siemssen} established a Banach contraction for the
semiclassical Einstein equation on FLRW spacetimes, using symmetry-reduced bounds
specific to cosmological backgrounds. Meda and Pinamonti~\cite{meda_pinamonti_stability}
showed that massive fields give polynomial decay of perturbations. The present Lemma
addresses the general (non-symmetric) case via the Hadamard decomposition.

\medskip

\begin{theorem}[Perturbative existence---massive fields --- \textbf{Sketch}]
\label{thm:banach}
Let $\bar{g}$ be a classical solution with sub-Planckian curvature
$R \ll \ell_P^{-2}$, \emph{admitting a compact Cauchy surface}, and let the matter
content consist of massive fields with
$m > 0$. Under the hypotheses of Lemma~\ref{lem:kernel_bound}, the self-consistency
map $\mathcal{F}$ is a contraction in a neighborhood of $\bar{g}$ in $H^s$, with
Lipschitz constant
\begin{equation}
\kappa \;=\; 8\pi G \cdot C_{\mathrm{inv}} \cdot
\sup_{g \in K_\rho}\bigl\|\mathcal{K}^{\mathrm{red}}[g]\bigr\|_{H^s \to H^{s-2}}
\;\leq\; C \cdot \frac{1}{2\pi}\!\left(\frac{m}{M_P}\right)^{\!2}\ell_P^2\, R
\;\ll\; 1
\label{eq:kappa}
\end{equation}
where $C_{\mathrm{inv}}$ is the norm of the inverse linearized Einstein operator
from~\emph{(A4)} and $C$ is a dimensionless constant. By the Banach fixed-point
theorem~\cite{banach}, there exists a unique self-consistent solution $g^*$ in a neighborhood of
$\bar{g}$, and the iteration $g_{n+1} = \mathcal{F}(g_n)$ converges exponentially:
$\|g_n - g^*\|_{H^s} \leq \kappa^n\|g_0 - g^*\|_{H^s}$.
\end{theorem}

\noindent\emph{Proof.} The Lipschitz constant of $\mathcal{F}$ is
\begin{equation}
\kappa = 8\pi G\,\bigl\|(G_{\mathrm{lin}})^{-1}\bigr\|_{H^{s-2}\to H^s}
\cdot \sup_{g \in K_\rho}
\bigl\|\mathcal{K}^{\mathrm{red}}[g]\bigr\|_{H^s \to H^{s-2}}.
\end{equation}
By assumption~(A4), the inverse linearized Einstein operator satisfies
$\|(G_{\mathrm{lin}})^{-1}\|_{H^{s-2}\to H^s} \leq C_{\mathrm{inv}}$, with
$C_{\mathrm{inv}} \sim 1/R$, since the harmonic gauge condition~(A4) renders the linearized Einstein operator elliptic with principal symbol determined by the background curvature scale~$R$; standard elliptic regularity estimates then give an inverse operator norm scaling as~$1/R$. By Lemma~\ref{lem:kernel_bound}, the response
kernel norm is bounded by $C_{\mathcal{K}} \cdot \hbar m^2/(4\pi)^2$. Combining:
\begin{equation}
\kappa \;\leq\; 8\pi G \cdot \frac{C_{\mathrm{inv}}C_{\mathcal{K}}\,\hbar\, m^2}{(4\pi)^2}
\;\sim\; \frac{\hbar G\, m^2}{2\pi}
\cdot \frac{1}{R}
\;=\; \frac{1}{2\pi}\!\left(\frac{m}{M_P}\right)^{\!2}\ell_P^2\, R.
\end{equation}
For sub-Planckian curvature ($R \ll \ell_P^{-2}$, i.e., $\ell_P^2 R \ll 1$)
and any known particle mass ($m \ll M_P$), we have $\kappa \ll 1$.
The Banach fixed-point theorem~\cite{banach} then guarantees a unique
fixed point in the closed ball around $\bar{g}$ of radius
$\rho/(1-\kappa)$, with geometric convergence.

\hfill$\square$

\smallskip\noindent\textbf{Gap (M1): dimensional inconsistency in the $\kappa$ chain.}
The proof asserts $C_{\mathrm{inv}} \sim 1/R$ (dimension $[\mathrm{length}]^2$,
consistent with an $H^{s-2}\to H^s$ operator norm), and then writes
$\kappa \sim (\hbar G m^2/2\pi)\cdot(1/R) = (1/2\pi)(m/M_P)^2\ell_P^2 R$. The two sides
of this equality are different functions of $R$ with different dimensions: with
$\hbar G = \ell_P^2$, the left side is $\ell_P^2 m^2/(2\pi R)$ (dimension
$[\mathrm{length}]^2$) while the right side is dimensionless; they coincide only if
$1/R = \ell_P^2 R$, i.e.\ the step silently sets $R = \ell_P^{-2}$. Followed literally,
the chain gives $\kappa \propto 1/R$, which \emph{diverges} in the flat-space limit
$R \to 0$ --- contraction would fail near Minkowski space, which is physically absurd
and contradicts the displayed formula $\kappa \propto R$. Which $R$-scaling (if either)
is correct is open; until the chain is repaired, the quantitative form of
eq.~\eqref{eq:kappa}, and every claim downstream of it (notably
Section~\ref{sec:planck}), is unsupported.

\emph{Resolution (Sketch).} The inconsistency is a hidden-reference-scale artifact. A
dimensionally homogeneous Sobolev norm carries a reference mass~$\mu$, $\|u\|_{H^t}^2 =
\langle u,(\mu^2-\Delta)^t u\rangle$, so an operator norm $\|A\|_{H^a\to H^b}$ carries an
explicit factor $\mu^{a-b}$ times the intrinsic dimension of~$A$. Then $C_{\mathrm{inv}} =
\|(G_{\mathrm{lin}})^{-1}\|_{H^{s-2}\to H^s}$ carries $\mu^{-2}$ (its ``$1/R$'' is really
$1/(\text{spectral gap of }G_{\mathrm{lin}})$, equal to $1/R$ only when $\mu^2\sim R$),
while $\|\mathcal{K}^{\mathrm{red}}\|_{H^s\to H^{s-2}}$ carries $\mu^{+2}$ (the asserted
$\hbar m^2/(4\pi)^2$ is licensed only when $\mu^2\sim m^2$). The displayed chain uses
\emph{two different} implicit scales --- $\mu^2=R$ in $C_{\mathrm{inv}}$ and $\mu^2=m^2$
in the kernel --- which is exactly why the unbalanced factor $\ell_P^2 R$ appears. A
\emph{single} explicit $\mu$ in both factors removes the inconsistency, at the cost of a
physical choice that collides with Gap~(M4):
\begin{itemize}
\item $\mu^2\sim m^2$ (the field's Compton scale): $\kappa\sim(m/M_P)^2$, \emph{independent
of $R$}, with a graceful flat-space limit ($\kappa$ unaffected as $R\to0$; no lower bound
on~$R$ is needed). \emph{This is schematic:} the Sobolev-index factors $\mu^{-2}$ (from
$C_{\mathrm{inv}}$) and $\mu^{+2}$ (from $\|\mathcal{K}^{\mathrm{red}}\|$) cancel, but
the surviving net physical scale depends on the intrinsic dimensions of the operator
product, which requires the joint (M1)$+$(M4) treatment to make precise. The corrected
Sketch-level bound below is the precise statement.
\item $\mu^2\sim R$: $\kappa\sim\ell_P^2 R$ (also $\to0$ as $R\to0$), but consistency then
forces the kernel scale to $\hbar R/(4\pi)^2$, not $\hbar m^2/(4\pi)^2$ --- precisely
Gap~(M4). The $(m/M_P)^2$ headline does not survive this reading.
\end{itemize}
The corrected Sketch-level statement is $\kappa \le C\,(m/M_P)^2\, f(\mu^2/m^2,\ell_P^2 R)$
with $f\to\mathrm{const}$ for $\mu^2\sim m^2$ and $f\sim R/m^2$ for $\mu^2\sim R$. The
theorem requires an explicit single~$\mu$, fixed jointly with~(M4) --- not a lower bound
on~$R$. The displayed $\kappa = (1/2\pi)(m/M_P)^2\ell_P^2 R$ is an artifact of mixing the
two scales and then setting $R=\ell_P^{-2}$; it is not a derivation. (This confirms the
Section~\ref{sec:planck} withdrawal from a second direction: on the defensible $\mu^2\sim
m^2$ reading $\kappa$ is $R$-independent, so no breakdown curvature is derived at all.)

\smallskip\noindent\textbf{Gap (M5): missing self-mapping step.}
The invocation of the Banach theorem on ``the closed ball around $\bar{g}$ of radius
$\rho/(1-\kappa)$'' requires, in addition to the contraction estimate, the displacement
bound $\|\mathcal{F}(\bar{g}) - \bar{g}\|_{H^s} \leq (1-\kappa)\rho$ --- i.e.\ that the
quantum correction evaluated \emph{at the background} is small in $H^s$. This is never
established; it is essentially an (A3)-type bound at $\bar{g}$. Contraction alone does
not yield existence without it. This is a genuine technical gap, subordinate to (M3): it
would still have to be discharged for any foliation-restricted repair (Section~\ref{sec:open}),
in the time-weighted norm that repair uses rather than in $H^s(\Sigma)$.

\medskip

For all known particles, $\kappa_{\text{mass}} \equiv (m/M_P)^2 \ll 1$:
\begin{center}
\begin{tabular}{lcc}
\toprule
Particle & Mass (kg) & $\kappa_{\text{mass}} = (m/M_P)^2$ \\
\midrule
Electron & $9.1\times10^{-31}$ & $\sim 10^{-45}$ \\
Proton & $1.7\times10^{-27}$ & $\sim 10^{-39}$ \\
Top quark & $3.1\times10^{-25}$ & $\sim 10^{-34}$ \\
\bottomrule
\end{tabular}
\end{center}

%======================================================================
\section{Non-perturbative existence via Schauder}
\label{sec:schauder}
%======================================================================

\subsection{Setup}

Let $(\mathcal{M},\bar{g})$ be a globally hyperbolic spacetime with \emph{compact} Cauchy
surface $\Sigma$. Fix $s > n/2 + 2 = 4$ and define the convex compact set:
\begin{equation}
K_\rho \equiv \bigl\{h \in H^s(\Sigma) : \|h-\bar{h}\|_{H^s} \leq \rho\bigr\}
\cap \mathcal{G}_{\text{harm}} \cap \mathcal{L}
\label{eq:K_rho}
\end{equation}
where $\mathcal{G}_{\text{harm}}$ is the affine set of metrics in linearized harmonic
gauge and $\mathcal{L}$ is the open set of Lorentzian-signature metrics. By the
Rellich--Kondrachov theorem on compact $\Sigma$, $K_\rho$ is compact in $H^{s'}$ for
any $s' < s$.

\subsection{Assumptions}

\begin{itemize}
\item[\textbf{(A1)}] \textit{Hadamard state existence.} For each $g\in K_\rho$, a
Hadamard state $|\Psi_g\rangle$ exists and $\langle\hat{T}_{\mu\nu}\rangle_{g,\text{ren}}$
is well-defined. \textbf{[Established:} Radzikowski~\cite{radzikowski},
Fulling--Sweeny--Wald~\cite{fulling_sweeny_wald}.\textbf{]}

\item[\textbf{(A2)}] \textit{Continuous metric dependence.} The map $g\mapsto
\langle\hat{T}_{\mu\nu}\rangle_{g,\text{ren}}$ is continuous $H^s\to H^{s-2}$.
\textbf{[Established:} Hollands--Wald~\cite{hollands_wald_wick}.\textbf{]}

\item[\textbf{(A3)}] \textit{Uniform stress-energy bound.} There exists $B(\rho)<\infty$
with $\|\langle\hat{T}_{\mu\nu}\rangle_{g,\text{ren}}\|_{H^{s-2}} \leq B(\rho)$ for all
$g\in K_\rho$. \textbf{[Conjecture.} Established for FRW spacetimes
\cite{pinamonti_siemssen}; open in general.\textbf{]}

\item[\textbf{(A4)}] \textit{Linearized Einstein well-posedness.} The linearized Einstein
equation in harmonic gauge has a unique solution $h\in H^s$ with
$\|h\|_{H^s} \leq C_{\text{inv}}\|S\|_{H^{s-2}}$.
\textbf{[Established:} Choquet-Bruhat~\cite{choquet_bruhat}.\textbf{]}

\item[\textbf{(A5)}] \textit{Order reduction.} Fourth-derivative terms from the trace
anomaly can be perturbatively eliminated. \textbf{[Established:} Parker--Simon
\cite{parker_simon}.\textbf{]}

\item[\textbf{(A6)}] \textit{Uniform contraction bound.}
(Included for completeness; used only in the Banach uniqueness argument of Section~\ref{sec:banach}, not in the Schauder existence argument below.)\newline
$\kappa_{\sup} \equiv 8\pi G\,C_{\text{inv}}\cdot\sup_{g\in K_\rho}\|\delta\langle\hat{T}_{\mu\nu}\rangle_{\text{ren}}/\delta g\|_{\text{op}} < 1$.
\textbf{[Sketch for massive fields, with a structural obstruction.}
Lemma~\ref{lem:kernel_bound} sketches the operator-norm bound. The displayed target
$\kappa \sim (m/M_P)^2\ell_P^2 R$ is not derived as written (Gap (M1): mixing two Sobolev
reference scales; the consistent single-scale form is $\kappa\sim(m/M_P)^2$, $R$-independent).
The governing gap (M3) is \emph{structural}: the response kernel is a Lorentzian hyperbolic
Green operator with causal-past support, and its reduction to a spatial operator on
$\Sigma$ requires a controlled global foliation beyond (A1)--(A6). The bound holds only on
such a (foliation-restricted) slicing; see Gaps (M1), (M3), (M4), (M5) in
Section~\ref{sec:banach}. Open for massless fields (the same (M3) gap binds).\textbf{]}
\end{itemize}

\subsection{Theorem}

\begin{theorem}[Non-perturbative existence (conditional on (A3))]
\label{thm:schauder}
Let $(\mathcal{M},\bar{g})$ be globally hyperbolic with compact Cauchy surface $\Sigma$,
and $\bar{g}$ a classical solution in linearized harmonic gauge. Under assumptions
\emph{(A1)--(A6)}, the self-consistency map $\mathcal{F}$ has at least one fixed point
$g^* \in K_\rho$ in the $H^{s'}$ topology for any $s' < s$.
\end{theorem}

\noindent\emph{Proof.} We verify the hypotheses of the Schauder fixed-point theorem
\cite{schauder}.

\smallskip\noindent\textit{Step 1: $K_\rho$ is convex and compact.} $K_\rho$ is the
intersection of a closed $H^s$ ball with the convex affine sets $\mathcal{G}_{\text{harm}}$
and $\mathcal{L}$, hence convex and closed in $H^s$. By Rellich--Kondrachov, it is compact
in $H^{s'}$.

\smallskip\noindent\textit{Step 2: $\mathcal{F}$ maps $K_\rho$ into $K_\rho$.} Given
$g\in K_\rho$: (A1) provides a Hadamard state with well-defined renormalized
stress-energy; (A5) reduces the fourth-order system to second order; (A4) gives
$\mathcal{F}(g) \in H^s$ with $\|\mathcal{F}(g)-\bar{h}\|_{H^s} \leq
C_{\text{inv}}\|S_{\mu\nu}\|_{H^{s-2}} \leq C_{\text{inv}}B(\rho)$ by (A3). Choose
$\rho = C_{\text{inv}}B(\rho)$ (a self-consistency condition on $\rho$ that is satisfiable
since $C_{\text{inv}}B(\rho)\to 0$ as $\rho\to 0$ for perturbative regimes). The harmonic
gauge and Lorentzian conditions are preserved by the linearized evolution.

\smallskip\noindent\textit{Step 3: $\mathcal{F}$ is continuous.} By (A2), $g\mapsto
S_{\mu\nu}[g]$ is continuous in $H^{s-2}$; by (A4), the Einstein solve is continuous;
composing, $\mathcal{F}: K_\rho\to K_\rho$ is continuous in $H^{s'}$.

\smallskip\noindent By the Schauder theorem~\cite{schauder}, $\mathcal{F}$ has a fixed
point in $K_\rho$.\hfill$\square$

\smallskip\noindent\textbf{Gap (M6): $K_\rho$ is not closed or convex as defined.}
Step~1 describes $\mathcal{L}$ as a ``convex affine set,'' but $\mathcal{L}$ is defined
in eq.~\eqref{eq:K_rho} as the \emph{open} set of Lorentzian-signature metrics, and the
set of Lorentzian metrics is neither affine nor convex (a convex combination of two
Lorentzian metrics need not be Lorentzian). As written, $K_\rho$ is therefore not closed,
and the compactness and convexity claimed in Step~1 do not follow. The expected repair is
cheap --- for $\rho$ small enough, Sobolev embedding keeps the entire closed $H^s$ ball
inside $\mathcal{L}$, so the intersection with $\mathcal{L}$ is vacuous --- but the proof
must be rewritten to say so; we flag rather than fix it here.

\smallskip\noindent\textbf{Gap (M7): the $\rho$-selection argument in Step~2 is incorrect.}
Step~2 claims the self-consistency condition $\rho = C_{\text{inv}}B(\rho)$ is satisfiable
``since $C_{\text{inv}}B(\rho)\to 0$ as $\rho\to 0$.'' This is false: as $\rho \to 0$,
$K_\rho \to \{\bar{g}\}$ and $B(\rho) \to \|\langle\hat{T}_{\mu\nu}\rangle_{\bar{g},
\text{ren}}\|_{H^{s-2}}$, which is the renormalized vacuum stress-energy of the background
and does not vanish. The correct route is different (e.g.\ $B$ is $O(\hbar)$-small and
approximately $\rho$-independent, so a fixed point of $\rho \mapsto C_{\text{inv}}B(\rho)$
exists for small $\hbar$), and the assertion that the harmonic-gauge and Lorentzian
conditions are ``preserved by the linearized evolution'' also requires argument (gauge
preservation needs source conservation, which holds with respect to $g$ but is applied in
a solve linearized around $\bar{g}$). Flagged, not fixed.

\noindent\emph{Remark.} Assumption~(A6) is not used in the Schauder argument (which
requires only continuity, not contractivity). It would be needed for a Banach fixed-point
argument guaranteeing uniqueness and constructive convergence --- an argument that is
itself presently a Sketch (Section~\ref{sec:banach}).

%======================================================================
\section{Effective initial-value formulation}
\label{sec:ivp}
%======================================================================

Theorem~\ref{thm:schauder} establishes global existence conditionally on~(A3).
Practical physics requires a
Cauchy formulation: given data on a spacelike hypersurface, what is the evolution? The
global fixed-point condition does not directly reduce to a Cauchy problem because the
Hadamard state is defined by a global microlocal condition. We derive an effective
formulation valid order-by-order in $\kappa$.

At order $\kappa^0$, the evolution is classical GR with given matter initial data.

At order $\kappa^1$, the semiclassical correction $\langle\hat{T}_{\mu\nu}\rangle^{(1)}$
is computed on the order-$\kappa^0$ background geometry. By Parker--Simon~\cite{parker_simon}
order reduction, the resulting system is second-order in time, with initial data for
$g_{\mu\nu}$ and $\partial_t g_{\mu\nu}$ on $\Sigma$ supplemented by the adiabatic vacuum
state at each order.

The order-reduced system at each order $\kappa^n$ is a second-order hyperbolic PDE,
well-posed by Choquet-Bruhat~\cite{choquet_bruhat}. The perturbative expansion is valid
while $\kappa \ll 1$; where that regime ends is unresolved pending the corrected
contraction estimate (Section~\ref{sec:planck}).

\noindent\textbf{Open problem.} Non-perturbative well-posedness of the order-reduced
semiclassical Einstein equation on general globally hyperbolic spacetimes without symmetry
assumptions has been conjectured by Ju\'arez-Aubry et al.~\cite{juarez_aubry} and
established for cosmological spacetimes~\cite{pinamonti_siemssen}. The general case
remains open.

%======================================================================
\section{Validity boundary: a withdrawn claim}
\label{sec:planck}
%======================================================================

An earlier revision of this paper claimed that the contraction
constant~\eqref{eq:kappa} satisfies $\kappa \to 1$ when $R \sim \ell_P^{-2}$, so that
the Planck scale emerges as the natural validity boundary of the framework --- derived
from the convergence requirement $\kappa < 1$ rather than imposed by hand. \textbf{This
claim is withdrawn.} It is contradicted by the paper's own formulas on either reading of
the $\kappa$ chain (Gap~(M1), Section~\ref{sec:banach}):

\begin{itemize}
\item Taking the displayed formula $\kappa = (1/2\pi)(m/M_P)^2\ell_P^2 R$ at face value,
at Planck curvature $R = \ell_P^{-2}$ one finds $\kappa = (m/M_P)^2/2\pi$, which by the
table in Section~\ref{sec:banach} is at most $\sim 10^{-34}$ for known particles ---
nowhere near $1$. On this reading, $\kappa \to 1$ only at curvatures
$\sim (M_P/m)^2$ \emph{above} the Planck scale.
\item Following the derivation chain literally ($C_{\mathrm{inv}} \sim 1/R$), one finds
$\kappa \propto 1/R$, which reaches $1$ at curvature $R \sim m^2$ --- the Compton scale
of the field, far \emph{below} the Planck scale --- and diverges in the flat-space limit.
\end{itemize}

\noindent Neither reading places the breakdown at the Planck scale, and the two disagree
with each other. Until the contraction estimate is repaired (Section~\ref{sec:open}),
the framework's validity boundary must be regarded as underived. The qualitative
expectation that semiclassical gravity fails at Planckian curvature remains physically
plausible, but it is not a result of this paper.

The framework also fails for observables that require summing over spacetime topologies
(e.g., the Page curve of black hole evaporation). Even when $\kappa \ll 1$---so the
perturbative expansion is valid---certain correlators receive $\mathcal{O}(1)$
contributions from topology-changing saddles absent from the fixed-topology
description. This is a limitation distinct from the (unresolved) convergence-boundary
question above, and is not self-diagnosable within the framework.

%======================================================================
\section{Open problems}
\label{sec:open}
%======================================================================

\begin{enumerate}

\item \textbf{The contraction estimate (Gaps M1, M3, M4, M5): a foliation-restricted
repair is the honest path.} The Banach argument of Section~\ref{sec:banach} is
permanently a Sketch at the stated generality. The gaps are not on equal footing
(analysis: \texttt{explorations/2026-06-16-FPE4-M3-structural-obstruction.md}).
\emph{(M3) is structural:} the non-local response kernel is a Green operator of the
Lorentzian hyperbolic operator, with causal-past support, and no reduction to a spatial
operator on a single Cauchy surface $\Sigma$ exists within (A1)--(A6) (sharing its root
with (M8) --- two faces of one obstruction: a genuine spatial reduction would simultaneously
close M8's well-definedness gap; the obstruction is mass-independent, so it is not cured
in the massless limit). The
only known repair replaces the single-slice $H^s(\Sigma)$ bound by a time-weighted energy
(Volterra) estimate on a finite causal slab, following Meda--Pinamonti--Siemssen
\cite{meda_pinamonti_siemssen,pinamonti_siemssen}; this requires a controlled global
foliation --- geometric input beyond (A1)--(A6) --- and yields existence ``Rigorous on a
[homogeneous / controlled] slicing,'' local in time, not the general statement of
Theorem~\ref{thm:banach}. \emph{(M1) is a hidden-reference-scale artifact} (diagnosed in Section~\ref{sec:banach};
corrected form schematic, net scale tied to the joint (M4) repair): with a single explicit
Sobolev scale $\mu$ the Sobolev-index mismatch is removed, giving $\kappa\sim(m/M_P)^2$
($R$-independent, schematic) for $\mu^2\sim m^2$, or $\kappa\sim\ell_P^2 R$ for
$\mu^2\sim R$ --- the latter forcing the local scale $\hbar R/(4\pi)^2$, i.e.\
\emph{(M4)}, so (M1) and (M4) are one joint repair. \emph{(M5)} (the self-mapping/displacement bound) remains a genuine technical gap
that any repair must still discharge, in the norm that repair uses. Restoring a Rigorous
label at the general stated scope is therefore not available; a foliation-restricted
re-promotion is.

\item \textbf{Well-definedness of the map $\mathcal{F}$ (M8, M9).}
Two structural gaps predate the demotion and remain open. \emph{(M8)}~Sections
\ref{sec:banach}--\ref{sec:schauder} work with fields in $H^s(\Sigma)$ on a single Cauchy
surface, while $\mathcal{F}$ is defined on spacetime geometries; the equivalence between
a spacetime fixed point and a fixed point of data on a chosen $\Sigma$ is never
constructed (the Hadamard condition is global), and the role of the chosen $\Sigma$ must
be shown to be pure gauge. \emph{(M9)}~The single-valuedness of $\mathcal{F}$ rests on
the assertion that the Hadamard state on the corrected geometry is ``the unique
perturbation'' of the background state; Hadamard states are not unique, and the
perturbation prescription that would make this precise is assumed, not constructed.
Relatedly, the ``Established'' tags on (A1) and (A2) are generous: Hollands--Wald
\cite{hollands_wald_wick} establishes local/analytic metric dependence of Wick
polynomials in the algebraic sense, not an $H^s \to H^{s-2}$ continuity statement, and
the deformation existence argument behind (A1) should be attributed and verified
precisely.

\item \textbf{Extension to massless fields.}
The binding obstruction for the massless extension is Gap~(M3), not an
independent infrared failure. Gap~(M3) is mass-independent: the non-local
response kernel is a Lorentzian Green operator with causal-past support
regardless of whether $m > 0$, and the algebraic-decay and Hilbert--Schmidt
arguments that appear in the massive Sketch are downstream of the spatial-operator
assumption Gap~(M3) already identifies as illegitimate. Closing the massless
extension therefore requires discharging (M3) first; any remaining IR issues
(e.g.\ algebraic kernel decay when $m = 0$) would be a secondary obstacle, not
the present bottleneck.

\item \textbf{Uniform stress-energy bound (A3).} The Pinamonti--Siemssen
result~\cite{pinamonti_siemssen} establishes this for FRW spacetimes. Extension to general
globally hyperbolic spacetimes with compact Cauchy surfaces requires controlling the
state-dependent part of $\langle\hat{T}_{\mu\nu}\rangle_{\text{ren}}$ uniformly over
$K_\rho$---an open problem.

\item \textbf{Relaxing the compact Cauchy surface assumption.} Both the Banach argument
(via Step~6 of Lemma~\ref{lem:kernel_bound} and the Hilbert--Schmidt bound) and the
Schauder argument require compactness of $K_\rho$, ensured by Rellich--Kondrachov on
compact $\Sigma$. Neither result therefore covers Minkowski space or asymptotically flat
backgrounds. Extension requires weighted Sobolev spaces and a different compactness
argument.

\item \textbf{Non-perturbative IVP well-posedness.} See Section~\ref{sec:ivp}.

\end{enumerate}

%======================================================================
\begin{thebibliography}{99}

\bibitem{oppenheim} J.~Oppenheim, A postquantum theory of classical gravity?,
\textit{Phys.\ Rev.\ X}\ \textbf{13}, 041040 (2023).

\bibitem{starobinsky} A.~A.~Starobinsky, A new type of isotropic cosmological models
without singularity, \textit{Phys.\ Lett.\ B}\ \textbf{91}, 99 (1980).

\bibitem{capper_duff} D.~M.~Capper and M.~J.~Duff, Trace anomalies in dimensional
regularization, \textit{Nuovo Cimento A}\ \textbf{23}, 173 (1974).

\bibitem{horowitz} G.~T.~Horowitz, Semiclassical relativity: The weak-field limit,
\textit{Phys.\ Rev.\ D}\ \textbf{21}, 1445 (1980).

\bibitem{phillips_hu} N.~G.~Phillips and B.~L.~Hu, Noise kernel in stochastic gravity and
stress energy bitensor, \textit{Phys.\ Rev.\ D}\ \textbf{63}, 104001 (2001).

\bibitem{hu_verdaguer} B.~L.~Hu and E.~Verdaguer, Stochastic gravity: Theory and
applications, \textit{Living Rev.\ Relativ.}\ \textbf{11}, 3 (2008).

\bibitem{radzikowski} M.~J.~Radzikowski, Micro-local approach to the Hadamard condition,
\textit{Commun.\ Math.\ Phys.}\ \textbf{179}, 529 (1996).

\bibitem{fulling_sweeny_wald} S.~A.~Fulling, M.~Sweeny, and R.~M.~Wald, Singularity
structure of the two-point function in quantum field theory in curved spacetime,
\textit{Commun.\ Math.\ Phys.}\ \textbf{63}, 257 (1978).

\bibitem{hollands_wald_wick} S.~Hollands and R.~M.~Wald, Local Wick polynomials and time
ordered products of quantum fields in curved spacetime, \textit{Commun.\ Math.\ Phys.}\
\textbf{223}, 289 (2001).

\bibitem{choquet_bruhat} Y.~Choquet-Bruhat, Th\'eor\`eme d'existence pour certains
syst\`emes d'\'equations aux d\'eriv\'ees partielles non lin\'eaires, \textit{Acta Math.}\
\textbf{88}, 141 (1952).

\bibitem{banach} S.~Banach, Sur les op\'erations dans les ensembles abstraits et leur
application aux \'equations int\'egrales, \textit{Fund.\ Math.}\ \textbf{3}, 133 (1922).

\bibitem{schauder} J.~Schauder, Der Fixpunktsatz in Funktionalr\"aumen,
\textit{Studia Math.}\ \textbf{2}, 171 (1930).

\bibitem{parker_simon} L.~Parker and J.~Z.~Simon, Einstein equation with quantum
corrections reduced to second order, \textit{Phys.\ Rev.\ D}\ \textbf{47}, 1339 (1993).

\bibitem{pinamonti_siemssen} N.~Pinamonti and D.~Siemssen, Global existence of solutions
of the semiclassical Einstein equation for cosmological spacetimes,
\textit{Commun.\ Math.\ Phys.}\ \textbf{334}, 171 (2015).

\bibitem{juarez_aubry} B.~A.~Ju\'arez-Aubry, B.~S.~Kay, T.~Miramontes, and D.~Sudarsky,
On the initial value problem for semiclassical gravity without and with quantum state
collapses, \textit{J.\ Cosmol.\ Astropart.\ Phys.}\ \textbf{2023}(01), 040 (2023).

\bibitem{meda_pinamonti_siemssen} P.~Meda, N.~Pinamonti, and D.~Siemssen,
Existence and uniqueness of solutions of the semiclassical Einstein equation in
cosmological models, \textit{Ann.\ Henri Poincar\'e}\ \textbf{22}, 3965 (2021).

\bibitem{meda_pinamonti_stability} P.~Meda and N.~Pinamonti,
Linear stability of semiclassical theories of gravity,
\textit{Ann.\ Henri Poincar\'e}\ \textbf{24}, 1211 (2023).

\end{thebibliography}

\end{document}
